% Generated from locale/tr/content/set-theory/z/infinity-again.tex; do not hand-edit.


\olfileid[tr]{sth}{z}{infinity-again}	
\olsection{Sonsuzluk}

Elimizdeki aksiyomlar, herhangi bir küme varsa sonsuz sayıda küme
bulunduğunu güvence altına almaya yeter. Bir kümenin var olduğunu varsayın;
o zaman $\emptyset$ vardır (\olref[sth][z][sep]{prop:emptyexists} gereği).
Şimdi her $x$ kümesi için $x \cup \{x\}$ kümesi
\olref[sth][z][pairs]{prop:pairsconsequences} gereği vardır. Bunu birkaç
kez uyguladığımızda şu kümeleri elde ederiz:
\begin{enumerate}
	\item[0.] $\emptyset$ %& $0$ & stage $0$\\
	\item[1.] $ \{\emptyset\}$ %& $1$ & stage $1$\\
	\item[2.] $\{\emptyset, \{\emptyset\}\}$ %& $2$ & stage $2$\\
	\item[3.] $\{\emptyset, \{\emptyset\}, \{\emptyset, \{\emptyset\}\}\}$ %&$3$ & stage $3$\\
	\item[4.] $\{\emptyset, \{\emptyset\}, \{\emptyset, \{\emptyset\}\}, \{\emptyset, \{\emptyset\}, \{\emptyset, \{\emptyset\}\}\}\}$% & $4$ & stage $4$ 
\end{enumerate}
Bu kümelerin birbirinden farklı olduğu denetlenebilir.

Numaralandırmaya $0$'dan başlamamızın birkaç nedeni var. Bunlardan biri
şudur: ``$n$'' diye etiketlediğimiz kümenin tam olarak $n$ elemanı olduğunu
ve (sezgisel olarak) $n$'inci aşamada oluşturulduğunu doğrulamak zor değildir.

Fakat bu bize \emph{sonsuz sayıda} küme verir; \emph{sonsuz bir kümenin},
yani sonsuz sayıda elemanı olan bir kümenin bulunduğunu güvence altına almaz.
Bu gerçekten önemlidir: Dedekind-sonsuz bir küme bulamazsak bir Dedekind
cebiri kuramayız. Oysa doğal sayılar kümesi olarak ele alabilmek için bir
Dedekind cebiri istiyoruz. (\olref[sfr][infinite][dedekindsproof]{sec}
ile karşılaştırın.)

Önemli olarak, şimdiye dek ortaya koyduğumuz aksiyomlar hiçbir sonsuz
kümenin varlığını güvence altına \emph{almaz}. Bu yüzden yeni bir aksiyom
ortaya koymalıyız:

\begin{axiom}[Sonsuzluk]
$\emptyset \in I$ olan ve $x \in I$ olduğunda $x \cup \{x\} \in I$ olan
bir $I$ kümesi vardır.
\begin{align*}
	\exists I( & (\exists o \in I)\forall x\ x \notin o \land {}\\
	& (\forall x \in I)(\exists s \in I)\forall z(z \in s \liff (z \in x \lor z = x)))
\end{align*}
\end{axiom}

Sonsuzluk Aksiyomu'nun bize verdiği $I$ kümesinin Dedekind-sonsuz olduğunu
görmek kolaydır. Ayırt edilmiş elemanı $\emptyset$, $I$ üzerindeki bire bir
fonksiyon ise $s(x)=x\cup\{x\}$'tir. Şimdi,
\olref[sfr][infinite][dedekind]{thm:DedekindInfiniteAlgebra},
Dedekind-sonsuz bir kümeden bir Dedekind cebirinin nasıl çıkarılacağını
gösterdi; biz de bunu doğal sayılar kümemiz olarak ele alacağız. Daha kesin
olarak:
\begin{defn}\ollabel{defnomega}
$I$, Sonsuzluk Aksiyomu'nun bize verdiği herhangi bir küme olsun. $s$,
$s(x)=x\cup\{x\}$ fonksiyonu olsun. $\omega =
\closureofunder{s}{\emptyset}$ olsun. $\omega$'nın elemanlarına
\emph{doğal sayılar} der ve $n$'nin, $s$'nin $\emptyset$'a $n$ kez
uygulanmasının sonucu olduğunu söyleriz.
\end{defn}

Şimdi birkaç paragraf önce ``$n$'' diye etiketlenen kümenin $n$ sayısı
\emph{olarak} ele alınacağını geriye dönüp doğrulayabilirsiniz.

Bu koşul koymanın önemini \olref[sth][z][nat]{sec} içinde tartışacağız.
Şimdilik bu, sezgisel bir sonucu ispatlamamızı sağlar:

\begin{prop}\ollabel{naturalnumbersarentinfinite}
Hiçbir doğal sayı Dedekind-sonsuz değildir.
\end{prop}

\begin{proof}
İspat tümevarımladır; yani
\olref[sfr][infinite][induction]{thm:dedinfiniteinduction} kullanılır.
Açıktır ki $0=\emptyset$ Dedekind-sonsuz değildir. Tümevarım adımı için
karşıt-tersi göstereceğiz: (saçma biçimde) $s(n)$ Dedekind-sonsuz ise $n$
de Dedekind-sonsuzdur.

Öyleyse $s(n)$'nin Dedekind-sonsuz olduğunu varsayın; yani
$\ran{f}\subsetneq\dom{f}=s(n)=n\cup\{n\}$ olan bir !!{injection} $f$
vardır. İki durum ele alınmalıdır.

\emph{Durum 1: $n \notin \ran{f}$.} O hâlde $\ran{f}\subseteq n$ ve
$f(n)\in n$'dir. $g=\funrestrictionto{f}{n}$ olsun; şimdi
$\ran{g}=\ran{f}\setminus\{f(n)\}\subsetneq n=\dom{g}$ olur. Dolayısıyla
$n$ Dedekind-sonsuzdur.

\emph{Durum 2: $n \in \ran{f}$.} Bir
$m\in\dom{f}\setminus\ran{f}$ sabitleyin ve tanım kümesi
$s(n)=n\cup\{n\}$ olan $h$ fonksiyonunu şöyle tanımlayın:
\[
h(x) =
	\begin{cases}
		f(x) & \text{$f(x) \neq n$ ise}\\
		m & \text{$f(x)=n$ ise.}
	\end{cases}
\]
Demek ki $h$ ile $f$, $h(f^{-1}(n))=m\neq n=f(f^{-1}(n))$ olan nokta
dışında her yerde aynıdır. $f$ bir !!a{injection} olduğundan
$n\notin\ran{h}$ ve $\ran{h}\subsetneq\dom{h}=s(n)$ olur. Şimdi Durum 1'in
gerekçesi kullanılarak $n$'nin Dedekind-sonsuz olduğu çıkar.
\end{proof}

Yine de Sonsuzluk Aksiyomu'nu nasıl \emph{gerekçelendireceğimiz} sorusu
kalır. Kısa yanıt, anlatmakta olduğumuz öyküye bir ilke daha eklememiz
gerektiğidir. Bu ilke şöyledir:
\begin{enumerate}
	\item[] \stagesinf. Sonsuz bir aşama vardır. Başka bir deyişle,
	(a) ilk aşama olmayan, (b) kendisinden önce bazı aşamalar bulunan, fakat
	(c) hemen önce gelen bir aşaması bulunmayan bir aşama vardır.
\end{enumerate}
Sonsuzluk Aksiyomu bu ilkeden doğrudan çıkar. Doğal sayı $n$'nin $n$'inci
aşamada oluşturulduğunu biliyoruz. Öyleyse $\omega$ kümesi ilk sonsuz
aşamada oluşturulur ve $\omega$'nın kendisi Sonsuzluk Aksiyomu'na tanıklık
eder.

Ancak bu bizi yalnızca \stagesinf{} ilkesini nasıl gerekçelendireceğimiz
sorusuna geri götürür. \stagessucc{} gibi bu da birikimli-yinelemeli
hiyerarşi hakkında anlattığımız öykünün açık bir parçası değildi. Üstelik
kümelerin aşama aşama oluşturulduğu yinelemeli hiyerarşi fikrinde, sürecin
\emph{sonsuz} bir aşama içerdiğini düşünmeye bizi zorlayan hiçbir şey
yoktur. Aşamaların doğal sayılar gibi sıralandığını düşünmek bütünüyle
tutarlı görünür.

Bu ise açık bir soruna yol açar. \olref[sfr][infinite][dedekindsproof]{sec}
içinde Dedekind'in (düşüncelerden oluşan) Dedekind-sonsuz bir kümenin
bulunduğuna ilişkin ``ispatını'' ele aldık. Bu size pek tatmin edici
gelmemiş olabilir. Fakat \stagesinf{} kümenin yinelemeli anlayışı (ya da
``düşünce yasaları'') tarafından ``bize zorunlu kılınmıyorsa'', bir
Dedekind-sonsuz kümenin bulunduğu savı için hâlâ içsel bir gerekçemiz yoktur.

Elbette bu konuda söylenecek çok daha fazla şey vardır. Ama umarız artık
bir aksiyomu (ya da ilkeyi) gerekçelendirmenin ne \emph{gerektirebileceği}
üzerine düşünmeye başlayacak bir noktadasınızdır. Bundan sonra
\stagesinf{} ilkesini olduğu gibi kabul edeceğiz.

